\chapter{Mesh, Materials, and Load Assembly: Inactive Paths}
\label{app:mesh}

This appendix chapter collects preprocessing routes that remain in the codebase
but are not part of the default parallel PMG benchmark path.

\section{Global Reference Assembly}

The repository still contains the more MATLAB-like route:
\begin{enumerate}
\item \code{assemble\_strain\_operator(...)},
\item \code{build\_elastic\_stiffness\_matrix(...)},
\item \code{vector\_volume(...)}.
\end{enumerate}

That path builds the global \code{B} matrix, global \code{WEIGHT}, global
\code{K\_elast}, and global \code{f\_V} exactly in the MATLAB style. It remains
important for validation, serial comparison, and direct formula tracing, but it
is not the normal PMG startup path.

\section{Asset-First Mesh Loading}

The active benchmark path uses canonical problem assets:
\begin{itemize}
\item \code{meshes/<asset>/definition.py} declares mesh variants, materials,
mechanics BCs, seepage BCs, hydraulic constants, and profiles,
\item \petscmod{assets/support/canonical\_gmsh.py} reads canonical linear Gmsh
meshes and promotes them to the requested element order,
\item \petscmod{problem\_asset\_runtime.py} resolves the selected
\code{asset}, \code{mesh\_variant}, and optional \code{profile}.
\end{itemize}

Older water-level and COMSOL-specific mesh loaders are historical migration
references only; seepage and coupled workflows now enter through the same asset
API as dry mechanics cases.

\section{Alternative Node Orderings}

\petscmod{mesh/reorder.py} still supports orderings beyond
\code{block\_metis}:
\begin{itemize}
\item \code{xyz},
\item \code{morton},
\item \code{rcm}.
\end{itemize}

Those orderings remain useful for experiments and diagnostics, but the committed
3D PMG benchmark is written around \code{block\_metis}.

\section{Why The Global Path Remains Important}

Although the owned-row path is mainline, the global path still serves three
purposes:
\begin{itemize}
\item it is the shortest conceptual bridge back to MATLAB,
\item it is the easiest correctness reference for distributed assembly,
\item it is the cleanest serial debugging path.
\end{itemize}
